% Study sheet: Volume I, Chapter 14, ML Operations.
% Numbers: \val{...} from ml_ops_studysheet_numbers.tex, written by ml_ops_studysheet_numbers.py.
\documentclass{studysheet}
\input{ml_ops_studysheet_numbers}
\sheetvolume{1}
\sheetchapter{14}{ML Operations}
\sheetsources{Numbers from the mlsysim registry and the chapter's ranking-service scenarios; assumed inputs are marked.}

\begin{document}
% ===================================================================== FRONT: LEARN
\frontside
\begin{bigidea}
A deployed model can be \hl{perfectly available and perfectly wrong}: every health check stays green
while the world drifts away from its training data. Operations keeps it right as a loop that never
ends, and each step of that loop is paid for in a \hl{machine resource}.
\end{bigidea}

\begin{leftcol}
\sheetsec{Key ideas}
\begin{ideas}
\idea{Available is not correct.}{PSI on inputs (investigate at \val{psiwarn}, escalate at \val{psicrit}) sees
  covariate drift but only warns. Labels confirm a loss and alone catch concept drift.}
\idea{Confirmation waits for labels.}{Time to confirm = labels needed $\div$ label arrival rate:
  \val{budget} labels take \val{minuteshigh} minutes at \val{ratehigh} per second, \val{dayslow} days at \val{ratelow} per day.}
\idea{Check parity per request.}{Replay recorded requests through both transforms and compare values by
  request ID. Matching histograms can hide wrong values.}
\idea{Diagnose before acting, data first.}{Pipeline, then skew, slices, time pattern, model. An abrupt
  drop points to a recent change; a gradual one to drift.}
\idea{Release safety costs memory.}{Each control holds a full second copy, whatever its traffic share.
  Rollback restores weights, feature schema and cache together.}
\idea{Maturity is the weakest suite.}{ML Test Score = the minimum of the data, model, infrastructure and
  monitoring scores. Raise the weakest first.}
\end{ideas}

\vspace{12pt}
\sheetsec{Technical debt at the boundaries}
\begin{sheettable}{@{}l L L@{}}
\thead{Pattern} & \thead{Symptom} & \thead{Guard}\\ \midrule
Boundary erosion & a change anywhere changes everything & modular interfaces\\
Correction cascade & an upstream fix breaks the models built on it & versioning; redesign, not patches\\
Undeclared consumer & an update silently breaks a hidden user & access control, interface contracts\\
Feedback loop & outputs shape future training data; decay hides & cohort monitoring, canaries\\
\end{sheettable}

\vspace{12pt}
\sheetsec{The running case: a ranking service}
DLRM's embedding tables are \val{tablegb} GB and an H100 holds \val{hbmgib} GiB (\val{hbmgb} GB), so one
copy of the model needs \val{percopy} H100s. Its node has \val{node}; a retraining run holds \val{runaccels}
of them for \val{runhours} hours.
\end{leftcol}%
\begin{rightcol}
\sheetsec{Formulas you will use}
\fcard{Accuracy decay}%
  {\AccuracyAfterDeployment(\DeploymentAge) = \InitialAccuracy\, e^{-\AccuracyDecayRate \DeploymentAge}}%
  {Accuracy fitted from outcomes falls at daily rate $\AccuracyDecayRate$ after a retraining run.}%
  {Fit from two readings: $\AccuracyDecayRate = \ln\!\big(\InitialAccuracy / \AccuracyAfterDeployment(\DeploymentAge)\big) / \DeploymentAge$.
   Holds for steady drift, not sudden breaks.}{}
\fcard{Optimal retraining interval}%
  {\RetrainingInterval^* \approx \sqrt{\frac{2\,\RetrainingRunCost}{\DailyQueryVolume\,\AccuracyPayoffPerQuery\,\InitialAccuracy\,\AccuracyDecayRate}}}%
  {Daily cost $= \RetrainingRunCost/\RetrainingInterval + \tfrac{1}{2}\DailyQueryVolume\AccuracyPayoffPerQuery\InitialAccuracy\AccuracyDecayRate\,\RetrainingInterval$:
   run cost $\RetrainingRunCost$ (compute, validation, deployment, risk) spread over $\RetrainingInterval$ days, plus
   average staleness. $\DailyQueryVolume$: queries a day; $\AccuracyPayoffPerQuery$: value per query of a whole unit of accuracy.}%
  {Valid while $\AccuracyDecayRate\RetrainingInterval$ is small. Four times any input moves $\RetrainingInterval^*$
   only twice. Retrain only when the recovered value exceeds $\RetrainingRunCost$.}{}

\vspace{10pt}
\sheetsec{Release controls}
\begin{sheettable}{@{}l L L L@{}}
\thead{Control} & \thead{Users exposed} & \thead{Extra memory} & \thead{What it buys}\\ \midrule
Shadow & none & full copy & output and latency check (scores every request twice)\\
Canary & small share & full copy & live quality evidence\\
A/B test & split & full copy & powered outcome comparison\\
Blue-green & all, after switch & full second environment & rollback by one route flip\\
Hot standby & none & full previous version & instant rollback\\
\end{sheettable}
\end{rightcol}

% ===================================================================== BACK: PRACTICE
\backside
\begin{answers}
\ans{Finish}{A: $\sqrt{\val{two} \times \text{\val{runa}} \div \text{\val{slope}}}$ = \val{tstara} days, barely moved: compute
  is only \val{computepct}\% of the run. B: $\sqrt{\val{two} \times \text{\val{runb}} \div \text{\val{slope}}}$ = \val{tstarb} days:
  cost fell \val{runratio}$\times$, so the interval fell $\sqrt{\val{runratio}}$ = \val{tratio}$\times$.}
\ans{1}{\val{budget} $\div$ \val{fraudrate} = \val{frauddays} days. Until then watch input PSI: a label-free
  warning that cannot show whether accuracy fell.}
\columnbreak
\ans{2}{No: \val{node} $-$ \val{percopy} $-$ \val{runaccels} = \val{freeaccels} free H100s, room for \val{freecopies}
  copy, which the hot standby needs. Start the canary when the run ends.}
\ans{3}{Data first: the drop is abrupt and follows a release. Replay those requests through both transforms by
  request ID; mismatched values mean a broken feature contract, not drift. Roll back the transform and its cache;
  retraining would learn the defect.}
\ans{4}{The minimum, \val{suitemin}; a model test leaves it unchanged. A feedback loop; cohort
  monitoring or a canary guards it.}
\end{answers}

\begin{leftcol}
\sheetsec{Worked example}
\begin{practicebox}\boxtitle{How often should the ranking service retrain?}
Accuracy is \val{accstart} at launch and \val{accweek} after \val{slidedays} days. $\DailyQueryVolume$ = \val{queries}
requests a day, $\AccuracyPayoffPerQuery$ = \$\val{value} per query, run cost $\RetrainingRunCost$ = \$\val{runcost}.\par\vspace{5pt}
{\centering
\begin{tikzpicture}[sheet]
  % x: 0..6 days -> 30..222 pt; y: 0..6000 dollars per day -> 0..62 pt
  \pgfmathsetmacro\sx{32} \pgfmathsetmacro\sy{62/6000}
  \foreach \d in {0,2,4,6}{\draw[gridline] ({30+\d*\sx},0) -- ({30+\d*\sx},62);
    \node[small, anchor=north] at ({30+\d*\sx},-3) {\d};}
  \foreach \c/\cl in {2000/2k,4000/4k,6000/6k}{\draw[gridline] (30,{\c*\sy}) -- (222,{\c*\sy});
    \node[small, anchor=east] at (28,{\c*\sy}) {\$\cl};}
  \node[small, anchor=north east] at (222,-12) {days between runs};
  \node[small, anchor=south west] at (0,66) {dollars per day};
  \draw[dim] (30,0) -- (222,0);
  \draw[muted, line width=0.8pt, domain=0.85:6, samples=60] plot ({30+\x*\sx},{\val{plotrun}/\x*\sy});
  \draw[muted, line width=0.8pt, dashed, domain=0:6] plot ({30+\x*\sx},{0.5*\val{plotslope}*\x*\sy});
  \draw[accent, line width=1.1pt, domain=1.0:6, samples=60] plot ({30+\x*\sx},{(\val{plotrun}/\x+0.5*\val{plotslope}*\x)*\sy});
  \pgfmathsetmacro\xs{30+\val{plottstar}*\sx}
  \pgfmathsetmacro\ys{(2*\val{plotrun}/\val{plottstar})*\sy}
  \draw[accent, dashed, line width=0.6pt] (\xs,0) -- (\xs,\ys);
  \filldraw[fill=accent, draw=accent] (\xs,\ys) circle (2.2pt);
  \node[note, anchor=south] at (\xs,{\ys+4}) {best: \val{tstar} days};
  \node[lab, text=muted, anchor=north] at ({30+5*\sx},{\val{plotrun}/5*\sy-2}) {run cost $\div$ days};
  \pgfmathsetmacro\ang{atan(0.5*\val{plotslope}*\sy/\sx)}
  \node[lab, anchor=south, rotate=\ang] at ({30+1.5*\sx},{0.5*\val{plotslope}*1.5*\sy+1}) {staleness};
  \node[lab, anchor=south east, text=accent] at (222,{(\val{plotrun}/6+0.5*\val{plotslope}*6)*\sy+2}) {total};
\end{tikzpicture}\par}\vspace{8pt}
\begin{steps}
\item Decay: $\AccuracyDecayRate$ = ln(\val{accstart} $\div$ \val{accweek}) $\div$ \val{slidedays} = \val{lnratio} $\div$ \val{slidedays}
  = \res{\val{gamma} per day}.
\item Staleness grows at \val{queries} $\times$ \val{value} $\times$ \val{accstart} $\times$ \val{gamma} = \res{\$\val{slope} per day$^2$}
  (daily cost adds half of it times $\RetrainingInterval$).
\item Interval: $\sqrt{2 \times \text{\val{runcost}} \div \text{\val{slope}}}$ = \res{\val{tstar} days}, where the two daily costs meet.
\end{steps}
\takeaway{The square root forgives rough inputs: four times the traffic only halves it, to \val{tstarfour} days.}
\end{practicebox}

\vspace{12pt}
\sheetsec{Your turn: finish it}
\begin{practicebox}\boxtitle{Faster training or faster validation?}
The run's compute is \val{runaccels} H100s $\times$ \val{runhours} h = \val{accelhours} accelerator-hours at
\$\val{price} per accelerator-hour, or \$\val{compute} of the \$\val{runcost} run. Same service, same slope.\par\vspace{4pt}
\textbf{Given:} (A) training twice as fast saves \$\val{halfsave}: $\RetrainingRunCost$ = \$\val{runa}.
(B) \assumed{} automated validation cuts $\RetrainingRunCost$ to \$\val{runb}.\par
\textbf{You:} find $\RetrainingInterval^*$ for each, and say which change moves the interval and why.
\end{practicebox}
\end{leftcol}%
\begin{rightcol}
\sheetsec{Common mistakes}
\mistake{``A \val{canarypct}\% canary costs \val{canarypct}\% of the hardware.''}%
  {Any request may look up any table row, so the canary holds all \val{tablegb} GB: the H100s holding
   tables go from \val{percopy} to \val{twocopies} while it serves one request in \val{canaryshare}.}{}
\mistake{``The training and serving histograms match, so there is no skew.''}%
  {Training gives requests A and B \val{sessiontrain} and \val{sessionserve} minutes; serving gives
   \val{sessionserve} and \val{sessiontrain}. Same histogram, both values wrong; only replay by request ID sees it.}{}
\mistake{``The drift alert fired, so retrain now.''}%
  {Divergence does not prove accuracy fell, and if a broken feature caused it, retraining bakes the defect
   in. Diagnose first; retrain only when the recovered value exceeds the run cost.}{}

\vspace{2pt}
\sheetsec{Check yourself}
\begin{qs}
\item \assumed{} a fraud model's policy needs \val{budget} labeled outcomes to confirm a drop, and chargebacks
  resolve \val{fraudrate} a day. How long until enough labels arrive? What can you watch meanwhile, and what
  can that not show?
\item On the ranking node a retraining run starts; rollback to the previous version must stay
  instant. Can a canary run during training? What do you schedule?
\item Click-through falls \val{ctrdrop}\% in an hour, just after a release of the serving feature transform;
  the weights did not change. Where do you look first, what separates a broken feature from drift, and
  should you retrain?
\item \assumed{} suite scores of data \val{suitedata}, model \val{suitemodel}, infrastructure \val{suiteinfra},
  monitoring \val{suitemon}. What is the ML Test Score? Would one more model test raise it? Its picks shape the clicks that train the next model: which debt pattern is this?
\end{qs}
\end{rightcol}
\end{document}
